\section{Parallel links, the exact-half moment, and a nested chain}
\label{sec:sharp-parallel-links}

The parallel-link argument is due to Fable 5.2.
Its bound depends only on the strict inequality above one half.
No positive lower bound on the excess is needed.

\begin{lemma}[Fable's parallel-link estimate]
\label{lem:endpoint-parallel-links}
Let $A,B$ be disjoint connected vertex sets in a connected multigraph $\Gamma$.
Let $P$ be a vertex set, and let $\xi$ be a uniform even subgraph.
Suppose its restriction outside $P$ is a forest with probability greater than one half.
For every integer $t\ge0$, at most $2^{t+1}$ components of $\Gamma-(A\cup B)$ satisfy all four conditions below:
\begin{enumerate}
\item They have an edge to each of $A,B$.
\item They avoid $P$.
\item They contain a cycle.
\item Their boundary has at most $t$ edges.
\end{enumerate}
\end{lemma}

\begin{proof}
If there is no component joining $A$ to $B$, the assertion is immediate.
Otherwise retain one such component $Q_0$.
The set $A\cup B\cup Q_0$ is connected.
For any two other complementary components $Q_i,Q_j$, the graph $\Gamma-Q_i-Q_j$ is connected.
It contains the retained connected set, and every other complementary component has an edge to that set.
No edge joins $Q_i$ and $Q_j$.

Thus the two regions have exactly one link.
The cycle-space pair law makes their zero-boundary events independent.
Conditional on both zero-boundary events, their internal even edge sets are independent and uniform.
This also follows from the direct product of the cycle spaces after deleting the two boundaries.

For each eligible component other than $Q_0$, let $E_i$ mean that its boundary is zero and its internal even edge set is nonzero.
The events $E_i$ are pairwise independent.
The zero-boundary condition imposes at most $t$ linear equations, so its probability is at least $2^{-t}$.
The cyclic internal graph has a nonzero even edge set with conditional probability at least one half.
Consequently $p_i:=\Pr(E_i)\ge2^{-t-1}$.

If there are $m$ such components, put $N=\sum_i\mathbf1_{E_i}$ and $\nu=\mathbb EN$.
Pairwise independence gives
\[
 \mathbb EN^2\le\nu^2+\nu,\qquad
 \Pr(N=0)\le\frac1{1+\nu},\qquad
 \nu\ge m2^{-t-1}.
\]
The forest event implies $N=0$.
Its probability exceeds one half, so $m<2^{t+1}$.
Including $Q_0$, if it was eligible, gives at most $2^{t+1}$ components.
\end{proof}


\begin{lemma}[A uniform link bound]
\label{lem:sharp-uniform-links}
Let $G$ be connected and cubic, and let $P$ be a nonempty connected marked set.
Assume that a uniform even subgraph is a forest outside $P$ with probability greater than one half.
Let $S$ be the union of a maximal vertex-disjoint family of cycles of length at most $D$ outside $P$.
Put $J=G-(P\cup S)$.
If $8\le D\le L$, $t\ge1$, and $2(t+2\lfloor D/8\rfloor)\le D$, then disjoint cycles $C,C'$ of $J$, each of length at most $L$, satisfy
\[
 \ell(C,C')\le1+\frac{2L}{t+1}+2^{t+1}+578(L/D)^2.
\]
\end{lemma}

\begin{proof}
Put $U=V(C)\cup V(C')$.
At most one component of $G-U$ meets $P$, because $P$ is connected.
At most $2L/(t+1)$ components have more than $t$ boundary edges: the total boundary of $U$ is at most $2L$.
Among the other components joining the two cycles, at most $2^{t+1}$ are cyclic by Lemma~\ref{lem:endpoint-parallel-links}.

Every remaining component $Y$ that joins both cycles is a tree, avoids $P$, and has boundary at most $t$.
The cubic degree identity gives $|Y|=|\delta Y|-2\le t-2$.
It contains no vertex of $S$.
Indeed, each chosen cycle avoids $U$, so any chosen cycle meeting $Y$ would lie entirely in that tree.
Thus $Y\subseteq J$.
Choose a path through $Y$ connecting $C$ to $C'$; its length is at most $t-1$.
Treat each direct edge between $C$ and $C'$ as another connecting path, of length one.
All these paths are internally disjoint.
Their endpoints are distinct, since each cycle vertex has at most one edge outside its cycle.

Put $a=\lfloor D/8\rfloor\ge D/16$.
Partition each cycle into at most $\lceil L/a\rceil\le17L/D$ consecutive arcs, each containing at most $a$ vertices.
Two connecting paths with endpoints in the same pair of arcs would give a cycle in $J$ of length at most
\[
 2t+2(a-1)\le2(t+2a)\le D.
\]
This contradicts the girth of $J$.
There is at most one path per pair of arcs, so the number of these links is at most $289(L/D)^2$.
The stated bound uses the harmless larger coefficient $578$.
Adding the four types of links proves the result.
\end{proof}

\begin{lemma}[The moment at a strict half threshold]
\label{lem:sharp-half-moment}
Use the preceding graph and marked set.
Let $\mathcal C$ be all cycles of $J$ of length at most $L$, and put
\[
 s=\lfloor(D-1)/2\rfloor,\qquad
 a_C=2^{-|C|},\qquad \mu=\sum_{C\in\mathcal C}a_C.
\]
Suppose every disjoint pair has at most $b$ links, and set
\[
 \eta=\frac32\,2^b L^3 2^{-s}.
\]
If $\mu>0$, the total weight $W_2$ of ordered disjoint pairs with exactly two links satisfies
\[
 W_2>\mu^2-(1+\eta)\mu.
\]
\end{lemma}

\begin{proof}
For each cycle use the variable $X_C$ from Section~\ref{sec:polynomial-foundations}, and put $Z=\sum_CX_C$.
Then $\mathbb EZ=\mu$.
The diagonal contribution to its second moment is at most $\mu$.
Intersecting distinct cycles contribute zero.
The pair law and the exceptional-partner estimate give
\[
 \mathbb EZ^2\le\mu^2+W_2+(1+\eta)\mu.
\]
Every cycle in $\mathcal C$ avoids $P$, so the forest event implies $Z=0$.
Its probability exceeds one half.
Consequently Cauchy--Schwarz gives
\[
 \mu^2\le\Pr(Z>0)\mathbb EZ^2<\frac12\mathbb EZ^2.
\]
Substitute the second-moment bound and rearrange.
\end{proof}

\begin{lemma}[A long nested chain]
\label{lem:sharp-nested-chain}
Under the preceding hypotheses, assume
\[
 \mu>20,\qquad \eta\le\frac12,\qquad L^2 2^{-s}\le\frac12.
\]
There are a cycle $C\subseteq J$ of length at most $L$, a neighbor $u$ of $C$, and a family of vertex-disjoint cycles $C_1,\ldots,C_T$ in $J$ such that
\[
 T\ge\frac{8\,2^s}{L^3}.
\]
All $C_i$ lie in the component of $G-V(C)$ containing $P$ and separate $u$ from $P$ there.
They admit an ordering for which their $u$-components $U_i$ in $G-V(C)-V(C_i)$ satisfy
\[
 U_i\subseteq U_{i+1},\qquad V(C_i)\subseteq U_{i+1}.
\]
\end{lemma}

\begin{proof}
For disjoint cycles $C,C'$, call $C'$ far behind $C$ if it lies outside the $P$-component of $G-V(C)$.
This relation is antisymmetric.
To see this, take a path from $P$ to its first vertex on $C$.
The path avoids the far component and hence avoids $C'$.
Connectedness of $C$ then puts all of $C$ in the $P$-component after deleting $C'$.
Thus ordered far pairs have total weight at most $\mu^2/2$.

For fixed $C$, cycles joined to it by an edge meet one of at most $L$ neighboring vertices.
The local weight bound makes their total weight at most $L^2 2^{-s}$.
Ordered adjacent pairs therefore have total weight at most $\mu L^2 2^{-s}$.
After excluding far and adjacent pairs, the preceding moment leaves ordered two-link pairs of total weight greater than
\[
 \frac12\mu^2-2\mu>8\mu.
\]
Weighted averaging gives a cycle $C$ whose remaining partners have total weight greater than eight.

Such a partner $C'$ has no edge to $C$ and has exactly two exterior linking components.
At least one avoids $P$ and contains a neighbor $u$ of $C$.
In $H=G-V(C)$, that partner separates $u$ from $P$.
There are at most $L$ choices of $u$.
For one choice the separating partners have total weight at least $8/L$.

Choose a maximal vertex-disjoint subfamily of these partners, of size $T$.
Its at most $TL$ vertices meet every partner in this family.
The local weight bound gives
\[
 \frac8L\le TL\cdot L2^{-s},
\]
which is the asserted lower bound.

Finally, disjoint connected separators of two connected sets have a consistent order.
Here is a direct justification in the present setting.
For one separator $A$, a path from $u$ to $P$ can be rerouted within $A$ between its first and last visits.
A disjoint connected separator $B$ must therefore lie in the $u$-component or the $P$-component of $H-A$.
If it lay in another component, that rerouted path would avoid $B$.
If $B$ lies on the $P$-side of $A$, then $A$ and the entire $u$-component of $H-A$ lie in the $u$-component of $H-B$.
The symmetric alternative reverses their order.
This orders every pair by inclusion of the $u$-components, giving the claimed nesting.
\end{proof}
